\subsection{Capture Region Analysis}\label{sec:capture_region}

We numerically compare the capture regions of two guidance strategies on a grid of initial conditions:
\begin{itemize}
    \item \textbf{VPP+LOS-rate:} the hierarchical policy that places a virtual pursuit point via PPO and tracks it with LOS-rate guidance.
    \item \textbf{PN direct:} the same LOS-rate guidance law, but with the virtual point fixed at the instantaneous target position (zero learned offset).
\end{itemize}

The grid spans initial range $R \in [500, 5000]$~m (10 points), heading error $\eta \in [-180^{\circ}, 180^{\circ}]$ (12 points), and speed ratio $v_o/v_t \in [0.8, 1.5]$ (4 points), for 480 initial conditions with 10 Monte Carlo episodes per condition. Figure~\ref{fig:capture_region_heatmap} replaces the previous table-only summary by showing both the VPP capture region and the VPP--PN success-rate difference across the grid.

\input{figures/capture_region_heatmap}

\subsubsection{Overall and sub-region comparison}

The pooled result remains nearly neutral: VPP+LOS-rate achieves 29.17\% mean success, PN direct achieves 28.75\%, and the overall difference is not statistically significant (Fisher exact test, $p = 0.669$). Sub-region tests also remain non-significant: high-ATA cases differ by +0.83 percentage points ($p=0.5242$), high-speed-ratio cases differ by +0.83 percentage points ($p=0.5372$), and the largest nominal advantage, high-ATA close/mid-range engagements, is +2.78 percentage points ($p=0.1930$). This is consistent with the physical intuition that, when the learned offset is small and the guidance gains are fixed to their default values, the VPP+LOS-rate law behaves similarly to classical proportional navigation.

\subsubsection{Physical interpretation}

The near-equivalence of VPP+LOS and PN direct in this study does \emph{not} imply that the VPP layer is unnecessary. Rather, it shows that under fixed, default parameters the hierarchical law can behave similarly to PN. The adaptive value of VPP arises only when the high-level policy is trained to choose context-dependent virtual-point offsets---lead pursuit in high-aspect engagements to avoid collision singularities, lag pursuit at high speed ratios to prevent overrun, and terminal blending near the minimum stand-off distance.

Because the capture-region grid freezes the policy and the gains, it measures the \emph{baseline kinematic capability} of the guidance law, not its \emph{learned adaptive capability}. The result therefore supports the interface-audit framing: VPP should be treated as a learnable augmentation of PN-like guidance, not as a replacement that is universally superior without additional evidence.

\subsubsection{Limitations}

The analysis uses the \texttt{simple} point-mass backend; JSBSim validation is future work. The PPO policy used in VPP+LOS was trained only for a short smoke run and is therefore not representative of a fully trained policy. The grid resolution (10~$\times$~12~$\times$~4) is coarse; finer grids and more episodes per cell are needed for high-confidence sub-region claims.
